%% ============================================================
%% Section : Empirical Strategy
%% ACV → Vote RN
%% ============================================================

\section{Empirical Strategy}
\label{sec:strategy}

\subsection{Identification Challenge}

Identifying the causal effect of ACV on RN vote shares faces two main obstacles. First, \textit{non-random selection}: ACV towns were chosen on the basis of observable indicators of urban decline---low median income, high unemployment, commercial vacancy---that are themselves correlated with prior RN support (Eatwell and Goodwin, 2018). A naive comparison of post-2018 RN trends in ACV versus non-ACV towns would conflate programme effects with the pre-existing differential trajectories of declining versus stable towns. Second, \textit{confounding temporal shocks}: the 2018--2019 window coincides with the \textit{Gilets Jaunes} movement (November 2018--April 2019), which may have independently reshaped political attitudes in towns experiencing economic hardship. Our strategy addresses both challenges.

\subsection{Propensity Score Matching}

As a first step, we reduce selection bias by restricting the control group to communes that closely resemble ACV towns in their pre-treatment socioeconomic profile. We estimate the probability of ACV designation---the propensity score $\hat{p}_i$---using a logistic regression of the treatment indicator on the following pre-treatment covariates (measured at or before 2017):
\begin{equation}
  \Pr(D_i = 1 \mid X_i) = \Lambda\left( \alpha + \beta_1 \text{Pop}_i + \beta_2 \text{RevMed}_i + \beta_3 \text{Chomage}_i + \beta_4 \text{Cadres}_i + \beta_5 \text{Ouvriers}_i + \beta_6 \bar{y}^{RN}_{i,\text{pre}} \right)
\end{equation}
where $\bar{y}^{RN}_{i,\text{pre}}$ is the commune's mean RN vote share across all pre-2018 electoral contests, and $\Lambda(\cdot)$ denotes the logistic cumulative distribution function.

We implement nearest-neighbour matching without replacement, retaining the three closest control communes for each treated commune (1:3 ratio) within a caliper of 0.25 standard deviations of the estimated propensity score \citep{Austin2011}. The AUC of the propensity score model is 0.813, indicating meaningful selection, and post-matching standardised mean differences (SMD) for all covariates fall below 0.2 (Table~\ref{tab:matching_balance}), satisfying the conventional threshold for covariate balance \citep{Stuart2010}.

\subsection{Difference-in-Differences Estimator}

We employ the difference-in-differences (DiD) estimator of \citet{Callaway2021} (hereafter CS2021) as our primary specification. Because all 222 ACV towns were designated simultaneously in March 2018, there is no staggered treatment timing in the strict sense, which means that the classical Two-Way Fixed Effects (TWFE) estimator does not suffer from the negative-weighting problem identified by \citet{Goodman-Bacon2021} in our main window. Nevertheless, we adopt CS2021 for three reasons: (i) it is robust to heterogeneous treatment effects across the commune distribution; (ii) it accommodates our pre-matched sample through the \texttt{xformla} argument; and (iii) its group-time ATT structure delivers transparent pre-trend tests.

Our estimating equation in the TWFE form---retained as a robustness benchmark---is:
\begin{equation}
  y_{it} = \alpha_i + \lambda_t + \delta \cdot (D_i \times \mathbf{1}[t \geq 2018]) + \varepsilon_{it}
  \label{eq:twfe}
\end{equation}
where $y_{it}$ is the RN vote share in commune $i$ at contest $t$; $\alpha_i$ is a commune fixed effect absorbing all time-invariant characteristics; $\lambda_t$ is an election-contest fixed effect capturing common temporal shocks; $D_i$ is the treatment indicator; and $\delta$ is the average treatment effect on the treated (ATT), interpreted as the change in RN vote share attributable to ACV designation.

Standard errors are clustered at the commune level throughout, to account for serial correlation in electoral outcomes within communes over time.

\subsection{Temporal Window and Pre-Trend Tests}

\paragraph{Main window: 2012--2019.} Our primary estimation window covers elections from 2012 to 2019. The upper bound is dictated by two post-2019 confounders: the \textit{Plan France Relance} (€100 billion, 2020--2022), which allocated large transfers to municipalities via channels correlated with urban decline, and the Covid-19 pandemic, which may have independently affected the political salience of state capacity in small and medium towns. Including post-2020 observations in the main specification would require modelling these additional treatments, substantially complicating the identification. We therefore treat the 2019 European elections as our single post-ACV outcome in the main specification and use 2022--2024 contests exclusively for robustness checks.

\paragraph{Event study.} To assess the plausibility of parallel trends, we estimate an event-study version of \eqref{eq:twfe} by interacting the treatment indicator with election-year dummies (normalised to zero at 2017, the last pre-treatment contest):
\begin{equation}
  y_{it} = \alpha_i + \lambda_t + \sum_{s \neq 2017} \gamma_s \cdot (D_i \times \mathbf{1}[t = s]) + \varepsilon_{it}
\end{equation}
The pre-treatment coefficients $\{\gamma_s\}_{s < 2018}$ should be jointly indistinguishable from zero under the null of parallel trends. We report these coefficients in Figure~\ref{fig:event_study}, along with 95\% confidence intervals based on commune-clustered standard errors.

\paragraph{Identification assumption.} The key identifying assumption is that, absent ACV designation, the electoral trajectory of ACV towns would have paralleled that of matched control towns. The pre-trend estimates (Section~\ref{sec:results}) are reassuringly small, but parallel trends is not directly testable. We therefore complement the pre-trend test with the sensitivity analysis of \citet{Rambachan2023}, which bounds the ATT under violations of parallel trends of bounded magnitude $\bar{M}$ (Section~\ref{sec:robustness}).

\subsection{Threats to Identification}

\paragraph{Selection into ACV.} The programme's selection was administrative rather than rule-based, precluding a regression discontinuity design. Our response is threefold: (i) propensity score matching to ensure covariate balance; (ii) pre-trend tests; and (iii) HonestDiD sensitivity bounds.

\paragraph{Gilets Jaunes.} The European elections of May 2019 took place two months after the main phase of the movement. If \textit{Gilets Jaunes} mobilisation was more intense in ACV towns (due to higher fuel price sensitivity in peri-urban areas), this could confound our post-treatment estimate. As a sensitivity check, we exclude the 10\% of communes with the highest \textit{Gilets Jaunes} density proxied by the number of local protest events (data from the Centre d'accès sécurisé aux données, when available), and compare the resulting estimate to the baseline.

\paragraph{Attrition and missing electoral data.} A small number of communes in our sample have missing data for one or more electoral contests, due to uncontested races or administrative irregularities. We treat these as missing-at-random and estimate the DiD on the unbalanced panel. Our robustness checks include a balanced subsample restricted to communes with complete data across all six contests.

\paragraph{Aggregation from polling station to commune.} ACV investment targets city centres, but voting outcomes are observed at the commune (municipality) level, which may include surrounding residential areas. This mismatch implies an attenuation bias in our estimates if programme effects are concentrated in central neighbourhoods. We note this as a limitation and suggest polling-station level analysis as a productive extension.
